% =====================================================================
% AAAI-2027 Reproducibility Checklist -- ANSWERED
%
% HOW TO USE: The authoritative template is the .tex from the AAAI-2027
% Author Kit (https://aaai.org/authorkit27/). It defines \question and the
% auto-numbering, and the instructions say to replace ONLY the
% "Type your response here" text. So: open the official checklist .tex and
% paste the one-word answer below under each matching question.
%
% This file is a self-contained, compilable convenience copy: it defines a
% minimal \question macro so you can build a filled PDF directly if you
% prefer. Answers are identical either way.
%
% NOTE on the three code-related items (2.8, 4.3, 4.4): they are answered
% "yes" on the assumption you upload the repository as the OpenReview
% "Code and Data Supplement". If you do NOT upload code with the submission,
% downgrade 2.8 to "no" and 4.3/4.4 to "partial".
% =====================================================================
\documentclass[letterpaper]{article}
\usepackage[margin=1in]{geometry}
\usepackage{enumitem}
\newcommand{\question}[2]{\vspace{6pt}\noindent\textbf{#1}~\textit{#2}\\}
\newcommand{\answer}[1]{\textbf{#1}}
\setlength{\parindent}{0pt}
\begin{document}

\begin{center}\Large\textbf{Reproducibility Checklist}\end{center}

\section*{1. General Paper Structure}

1.1.~\question{Includes a conceptual outline and/or pseudocode description of AI
methods introduced}{(yes/partial/no/NA)}
\answer{partial}

1.2.~\question{Clearly delineates statements that are opinions, hypothesis, and
speculation from objective facts and results}{(yes/no)}
\answer{yes}

1.3.~\question{Provides well-marked pedagogical references for less-familiar
readers to gain background necessary to replicate the paper}{(yes/no)}
\answer{yes}

\section*{2. Theoretical Contributions}

2.1.~\question{Does this paper make theoretical contributions?}{(yes/no)}
\answer{yes}

2.2.~\question{All assumptions and restrictions are stated clearly and
formally}{(yes/partial/no)}
\answer{yes}

2.3.~\question{All novel claims are stated formally (e.g., in theorem
statements)}{(yes/partial/no)}
\answer{yes}

2.4.~\question{Proofs of all novel claims are included}{(yes/partial/no)}
\answer{yes}

2.5.~\question{Proof sketches or intuitions are given for complex and/or novel
results}{(yes/partial/no)}
\answer{yes}

2.6.~\question{Appropriate citations to theoretical tools used are
given}{(yes/partial/no)}
\answer{yes}

2.7.~\question{All theoretical claims are demonstrated empirically to
hold}{(yes/partial/no/NA)}
\answer{yes}

2.8.~\question{All experimental code used to eliminate or disprove claims is
included}{(yes/no/NA)}
\answer{yes}

\section*{3. Dataset Usage}

3.1.~\question{Does this paper rely on one or more datasets?}{(yes/no)}
\answer{yes}

3.2.~\question{A motivation is given for why the experiments are conducted on the
selected datasets}{(yes/partial/no/NA)}
\answer{yes}

3.3.~\question{All novel datasets introduced in this paper are included in a data
appendix}{(yes/partial/no/NA)}
\answer{NA}

3.4.~\question{All novel datasets introduced in this paper will be made publicly
available upon publication of the paper with a license that allows free usage for
research purposes}{(yes/partial/no/NA)}
\answer{NA}

3.5.~\question{All datasets drawn from the existing literature are accompanied by
appropriate citations}{(yes/no/NA)}
\answer{yes}

3.6.~\question{All datasets drawn from the existing literature are publicly
available}{(yes/partial/no/NA)}
\answer{yes}

3.7.~\question{All datasets that are not publicly available are described in
detail, with explanation why publicly available alternatives are not
scientifically satisficing}{(yes/partial/no/NA)}
\answer{NA}

\section*{4. Computational Experiments}

4.1.~\question{Does this paper include computational experiments?}{(yes/no)}
\answer{yes}

4.2.~\question{This paper states the number and range of values tried per
(hyper-)parameter during development of the paper, along with the criterion used
for selecting the final parameter setting}{(yes/partial/no/NA)}
\answer{partial}

4.3.~\question{Any code required for pre-processing data is included in the
appendix}{(yes/partial/no)}
\answer{yes}

4.4.~\question{All source code required for conducting and analyzing the
experiments is included in a code appendix}{(yes/partial/no)}
\answer{yes}

4.5.~\question{All source code required for conducting and analyzing the
experiments will be made publicly available upon publication of the paper with a
license that allows free usage for research purposes}{(yes/partial/no)}
\answer{yes}

4.6.~\question{All source code implementing new methods have comments detailing
the implementation, with references to the paper where each step comes
from}{(yes/partial/no)}
\answer{partial}

4.7.~\question{If an algorithm depends on randomness, then the method used for
setting seeds is described in a way sufficient to allow replication of
results}{(yes/partial/no/NA)}
\answer{yes}

4.8.~\question{This paper specifies the computing infrastructure used for running
experiments (hardware and software), including GPU/CPU models; amount of memory;
operating system; names and versions of relevant software libraries and
frameworks}{(yes/partial/no)}
\answer{partial}

4.9.~\question{This paper formally describes evaluation metrics used and explains
the motivation for choosing these metrics}{(yes/partial/no)}
\answer{yes}

4.10.~\question{This paper states the number of algorithm runs used to compute
each reported result}{(yes/no)}
\answer{yes}

4.11.~\question{Analysis of experiments goes beyond single-dimensional summaries
of performance (e.g., average; median) to include measures of variation,
confidence, or other distributional information}{(yes/no)}
\answer{yes}

4.12.~\question{The significance of any improvement or decrease in performance is
judged using appropriate statistical tests (e.g., Wilcoxon
signed-rank)}{(yes/partial/no)}
\answer{partial}

4.13.~\question{This paper lists all final (hyper-)parameters used for each
model/algorithm in the paper's experiments}{(yes/partial/no/NA)}
\answer{yes}

\end{document}
